% ==============================================================================
% SLIDES - SEMANA 03: ALINHAMENTO ESTRATÉGICO DE TI & PETI
% ==============================================================================

% 1. CAPA
\senaicover
  {Alinhamento Estratégico \& PETI}
  {Modelos de Alinhamento (SAM), Plano Estratégico de TI e Matriz SWOT}
  {Unidade Curricular: Governança de TI}
  {Prof. André Souza}

% 2. AGENDA
\begin{senaiframe}{Visão Geral \& Agenda}{Estrutura da Aula}
  \begin{columns}[t]
    \begin{column}{0.48\textwidth}
      \begin{tcolorbox}[senaibluecard, title={1. Modelo SAM}]
        \begin{itemize}
          \item Alinhamento entre Negócio e Tecnologia (Luftman).
          \item As 4 perspectivas de integração estratégica.
        \end{itemize}
      \end{tcolorbox}
      \vspace{4pt}
      \begin{tcolorbox}[senaiambercard, title={3. Diagnóstico Situacional}]
        \begin{itemize}
          \item Mapeamento de dores e gargalos operacionais.
          \item Matriz SWOT aplicada à TI (FOFA).
        \end{itemize}
      \end{tcolorbox}
    \end{column}
    \begin{column}{0.48\textwidth}
      \begin{tcolorbox}[senairedcard, title={2. O PETI}]
        \begin{itemize}
          \item Estrutura e fases de elaboração do PETI.
          \item Vinculação com o Plano Estratégico do Negócio.
        \end{itemize}
      \end{tcolorbox}
      \vspace{4pt}
      \begin{tcolorbox}[senaigreencard, title={4. Capítulo 2 do PDGTI}]
        \begin{itemize}
          \item Diagnóstico da Situação Atual da TI.
          \item Avaliação preliminar de maturidade dos processos.
        \end{itemize}
      \end{tcolorbox}
    \end{column}
  \end{columns}
\end{senaiframe}

% 3. DESAFIO DO ALINHAMENTO
\begin{senaiframe}{Teoria • Módulo 1}{O Desafio do Alinhamento TI-Negócio}
  \begin{columns}[t]
    \begin{column}{0.48\textwidth}
      \begin{tcolorbox}[senairedcard, title={O Abismo Tradicional}]
        \begin{itemize}
          \item Negócio fala a língua de mercado e faturamento.
          \item TI fala a língua de servidores, código e suporte.
          \item Projetos entregues que não resolvem necessidades reais.
        \end{itemize}
      \end{tcolorbox}
    \end{column}
    \begin{column}{0.48\textwidth}
      \begin{tcolorbox}[senaigreencard, title={A TI Alinhada}]
        \begin{itemize}
          \item Projetos de tecnologia guiados pelas metas da empresa.
          \item Linguagem unificada em métricas de negócio.
        \end{itemize}
      \end{tcolorbox}
    \end{column}
  \end{columns}
\end{senaiframe}

% 4. O MODELO SAM
\begin{senaiframe}{Teoria • Módulo 1}{O Modelo de Alinhamento Estratégico (SAM)}
  \begin{columns}[t]
    \begin{column}{0.48\textwidth}
      \begin{tcolorbox}[senaibluecard, title={Estratégia vs. Operação}]
        \begin{itemize}
          \item Harmonia entre o ambiente Externo (Mercado) e Interno (Operação).
          \item Intersecção entre Estratégia de Negócio e Estratégia de TI.
        \end{itemize}
      \end{tcolorbox}
    \end{column}
    \begin{column}{0.48\textwidth}
      \begin{tcolorbox}[senaicard, title={Henderson \& Venkatraman}]
        \begin{itemize}
          \item A TI não pode ser tratada de forma isolada.
          \item Mudanças no mercado exigem adequação na infraestrutura.
        \end{itemize}
      \end{tcolorbox}
    \end{column}
  \end{columns}
\end{senaiframe}

% 5. PERSPECTIVAS DE LUFTMAN
\begin{senaiframe}{Teoria • Módulo 1}{As 4 Perspectivas de Alinhamento (Luftman)}
  \begin{columns}[t]
    \begin{column}{0.48\textwidth}
      \begin{tcolorbox}[senaibluecard, title={Execução \& Transformação}]
        \begin{itemize}
          \item \textbf{Execução Estratégica:} O negócio impulsiona a TI.
          \item \textbf{Transformação Tecnológica:} TI cria novas ofertas.
        \end{itemize}
      \end{tcolorbox}
    \end{column}
    \begin{column}{0.48\textwidth}
      \begin{tcolorbox}[senaiambercard, title={Competição \& Serviços}]
        \begin{itemize}
          \item \textbf{Potencial Competitivo:} Inovação disruptiva.
          \item \textbf{Nível de Serviço:} Eficiência operacional interna.
        \end{itemize}
      \end{tcolorbox}
    \end{column}
  \end{columns}
\end{senaiframe}

% 6. O QUE É O PETI
\begin{senaiframe}{Teoria • Módulo 2}{O Plano Estratégico de TI (PETI)}
  \begin{columns}[t]
    \begin{column}{0.48\textwidth}
      \begin{tcolorbox}[senaibluecard, title={Conceito do PETI}]
        \begin{itemize}
          \item Documento formal de diretrizes tecnológicas de médio/longo prazo.
          \item Vinculado diretamente ao Plano Estratégico do Negócio (PEB).
        \end{itemize}
      \end{tcolorbox}
    \end{column}
    \begin{column}{0.48\textwidth}
      \begin{tcolorbox}[senaicard, title={Objetivos Centrais}]
        \begin{itemize}
          \item Direcionar investimentos futuros de TI.
          \item Evitar contratações de sistemas incompatíveis.
        \end{itemize}
      \end{tcolorbox}
    \end{column}
  \end{columns}
\end{senaiframe}

% 7. FASES DO PETI - 1 E 2
\begin{senaiframe}{Teoria • Módulo 2}{Fases de Elaboração do PETI (Fases 1 \& 2)}
  \begin{columns}[t]
    \begin{column}{0.48\textwidth}
      \begin{tcolorbox}[senaibluecard, title={Fase 1: Contexto do Negócio}]
        \begin{itemize}
          \item Entendimento das metas corporativas (expansão, receita).
          \item Mapeamento das prioridades estratégicas da diretoria.
        \end{itemize}
      \end{tcolorbox}
    \end{column}
    \begin{column}{0.48\textwidth}
      \begin{tcolorbox}[senairedcard, title={Fase 2: Diagnóstico da TI}]
        \begin{itemize}
          \item Levantamento da arquitetura atual (As-Is).
          \item Mapeamento de problemas de infraestrutura e sistemas.
        \end{itemize}
      \end{tcolorbox}
    \end{column}
  \end{columns}
\end{senaiframe}

% 8. FASES DO PETI - 3 E 4
\begin{senaiframe}{Teoria • Módulo 2}{Fases de Elaboração do PETI (Fases 3 \& 4)}
  \begin{columns}[t]
    \begin{column}{0.48\textwidth}
      \begin{tcolorbox}[senaiambercard, title={Fase 3: Arquitetura Futura}]
        \begin{itemize}
          \item Definição do estado desejado (To-Be).
          \item Escolha de tecnologias, nuvem e padrões de segurança.
        \end{itemize}
      \end{tcolorbox}
    \end{column}
    \begin{column}{0.48\textwidth}
      \begin{tcolorbox}[senaigreencard, title={Fase 4: Plano de Transição}]
        \begin{itemize}
          \item Roadmap de projetos, orçamentos e prazos.
          \item Definição dos responsáveis por cada iniciativa.
        \end{itemize}
      \end{tcolorbox}
    \end{column}
  \end{columns}
\end{senaiframe}

% 9. DIAGNÓSTICO DA SITUAÇÃO ATUAL
\begin{senaiframe}{Teoria • Módulo 3}{O Diagnóstico Situacional da TI (As-Is)}
  \begin{columns}[t]
    \begin{column}{0.48\textwidth}
      \begin{tcolorbox}[senaicard, title={Levantamento de Dores}]
        \begin{itemize}
          \item Identificação de lentidão, indisponibilidades e retrabalho.
          \item Entrevistas com gestores das áreas de negócio.
        \end{itemize}
      \end{tcolorbox}
    \end{column}
    \begin{column}{0.48\textwidth}
      \begin{tcolorbox}[senairedcard, title={Gargalos de Integração}]
        \begin{itemize}
          \item Ilhas de informação e planilhas paralelas de controle.
          \item Falta de automação em processos operacionais.
        \end{itemize}
      \end{tcolorbox}
    \end{column}
  \end{columns}
\end{senaiframe}

% 10. ESTRUTURA DA MATRIZ SWOT
\begin{senaiframe}{Teoria • Módulo 3}{A Matriz SWOT Aplicada à TI}
  \begin{columns}[t]
    \begin{column}{0.48\textwidth}
      \begin{tcolorbox}[senaibluecard, title={Fatores Internos (Sob Controle)}]
        \begin{itemize}
          \item \textbf{Forças (Strengths):} Recursos e capacidades positivas da TI.
          \item \textbf{Fraquezas (Weaknesses):} Deficiências internas de TI.
        \end{itemize}
      \end{tcolorbox}
    \end{column}
    \begin{column}{0.48\textwidth}
      \begin{tcolorbox}[senaiambercard, title={Fatores Externos (Fora do Controle)}]
        \begin{itemize}
          \item \textbf{Oportunidades (Opportunities):} Tendências de mercado.
          \item \textbf{Ameaças (Threats):} Riscos externos e regulação.
        \end{itemize}
      \end{tcolorbox}
    \end{column}
  \end{columns}
\end{senaiframe}

% 11. QUADRANTES INTERNOS (S & W)
\begin{senaiframe}{Teoria • Módulo 3}{SWOT: Forças (S) \& Fraquezas (W) na TI}
  \begin{columns}[t]
    \begin{column}{0.48\textwidth}
      \begin{tcolorbox}[senaigreencard, title={Exemplos de Forças (S)}]
        \begin{itemize}
          \item Equipe qualificada em metodologias ágeis.
          \item Servidores redundantes e alta disponibilidade.
        \end{itemize}
      \end{tcolorbox}
    \end{column}
    \begin{column}{0.48\textwidth}
      \begin{tcolorbox}[senairedcard, title={Exemplos de Fraquezas (W)}]
        \begin{itemize}
          \item Ausência de plano de disaster recovery.
          \item Dependência de pessoas-chave (key-person risk).
        \end{itemize}
      \end{tcolorbox}
    \end{column}
  \end{columns}
\end{senaiframe}

% 12. QUADRANTES EXTERNOS (O & T)
\begin{senaiframe}{Teoria • Módulo 3}{SWOT: Oportunidades (O) \& Ameaças (T) na TI}
  \begin{columns}[t]
    \begin{column}{0.48\textwidth}
      \begin{tcolorbox}[senaibluecard, title={Exemplos de Oportunidades (O)}]
        \begin{itemize}
          \item Adoção de Inteligência Artificial para automação.
          \item Migração para serviços gerenciados SaaS.
        \end{itemize}
      \end{tcolorbox}
    \end{column}
    \begin{column}{0.48\textwidth}
      \begin{tcolorbox}[senaiambercard, title={Exemplos de Ameaças (T)}]
        \begin{itemize}
          \item Ataques cibernéticos e Ransomware.
          \item Exigências regulatórias rigorosas (LGPD).
        \end{itemize}
      \end{tcolorbox}
    \end{column}
  \end{columns}
\end{senaiframe}

% 13. ESTUDO DE CASO PRÁTICO
\begin{senaiframe}{Estudo de Caso}{Grupo Hospitalar: Reorientação do PETI}
  \begin{columns}[t]
    \begin{column}{0.48\textwidth}
      \begin{tcolorbox}[senairedcard, title={Desalinhamento Constatado}]
        \begin{itemize}
          \item Meta da empresa: reduzir tempo na emergência de 45 para 15 min.
          \item Foco da TI: migração de servidores do setor financeiro.
        \end{itemize}
      \end{tcolorbox}
    \end{column}
    \begin{column}{0.48\textwidth}
      \begin{tcolorbox}[senaigreencard, title={Solução via PETI}]
        \begin{itemize}
          \item Cancelamento de projetos secundários.
          \item Prioridade: Wi-Fi na emergência e prontuário em tablets -> 12 min.
        \end{itemize}
      \end{tcolorbox}
    \end{column}
  \end{columns}
\end{senaiframe}

% 14. REFERÊNCIAS BIBLIOGRÁFICAS
\begin{senaiframe}{Referências Bibliográficas}{Fontes \& Leitura Recomendada}
  \vspace{0.5cm}
  \begin{tcolorbox}[senaicard, title={Obras de Referência}]
    \begin{itemize}
      \item \textbf{FERNANDES, Aguinaldo Aragon; ABREU, Vladimir Ferraz.} *Implantando a governança de TI: da estratégia à gestão de processos e serviços.* 4. ed. Rio de Janeiro: Brasport, 2014.
      \item \textbf{LUFTMAN, Jerry.} *Managing the Information Technology Resource: Leadership in the Digital Age.* Pearson, 2003.
      \item \textbf{HENDERSON, J. C.; VENKATRAMAN, N.} *Strategic Alignment: Leveraging Information Technology for Transforming Organizations.* IBM Systems Journal, 1993.
    \end{itemize}
  \end{tcolorbox}
\end{senaiframe}

% 15. APLICAÇÃO NO PDGTI
\begin{senaiframe}{Aplicação no PDGTI}{Desenvolvimento do Capítulo 2}
  \vspace{0.5cm}
  \begin{tcolorbox}[senaigreencard, title={Capítulo 2 – Diagnóstico da Situação Atual}]
    \begin{itemize}
      \item \textbf{2.1 Matriz SWOT da TI:} Mapeamento detalhado dos 4 quadrantes.
      \item \textbf{2.2 Dores e Indisponibilidades:} Levantamento de falhas operacionais.
      \item \textbf{2.3 Avaliação do Desalinhamento:} Impacto das falhas no negócio.
      \item \textbf{2.4 Avaliação de Maturidade:} Nota de 1 a 5 para os processos.
    \end{itemize}
  \end{tcolorbox}
\end{senaiframe}
